\chapter{LITERATURE REVIEW}

Solar power generation is inherently variable because photovoltaic output depends on fluctuating meteorological conditions such as solar irradiance, ambient temperature, humidity, wind speed, atmospheric pressure, and cloud-related phenomena. These weather variables affect the amount of solar energy received by photovoltaic panels and influence the stability of solar power output. Previous studies have shown that solar radiation, sunshine duration, Julian day, cloud cover, and temporal meteorological variables are important predictors in photovoltaic generation and solar irradiance forecasting \parencite{CarvalhoJunior2025, Pongphaw2025, TsaiLo2025}. This variability introduces significant forecasting challenges because accurate predictions are needed for grid stability, energy scheduling, and storage management. Machine learning has emerged as an important approach for addressing these challenges because it can capture nonlinear relationships between weather sensor variables and solar power generation. However, the existing literature does not present a universally superior forecasting model because forecasting performance is influenced by model selection, weather sensor quality, preprocessing procedures, feature engineering, temporal resolution, forecasting horizon, and dataset context \parencite{JnanaVarshitha2025, Bakht2025, Eren2025, Alhassan2025}.

The objective of this study, Machine Learning-Based Forecasting of Solar Power Generation Using Weather Sensor Data, is to establish an effective ML framework for solar power forecasting that utilises weather sensor data while addressing gaps identified in previous research. This literature review examines the role of weather sensor variables, compares machine learning models used in solar forecasting, evaluates the importance of preprocessing and feature engineering, and discusses interpretability and generalizability in model evaluation. By doing so, this review directly supports the research gap concerning the difficulty of selecting and justifying the most suitable forecasting model for a specific weather-sensor-based solar power forecasting task \parencite{JnanaVarshitha2025, Alhassan2025, ArellanosTafur2025}.

Existing literature shows that although machine learning models can predict solar power effectively, their performance is often sensitive to dataset characteristics, geographical conditions, weather variables, sensor quality, temporal resolution, preprocessing procedures, and forecasting horizon. This variability highlights the need for a systematic comparison of models and feature sets under a consistent evaluation procedure \parencite{ArellanosTafur2025, Alhassan2025, CarvalhoJunior2025}. In addition, interpretability methods are needed to ensure that the selected model is not only accurate but also able to provide insight into the weather variables that influence solar power prediction. SHAP- and LIME-based studies show that feature influence analysis can improve understanding of how weather-based and time-based variables contribute to model predictions \parencite{JnanaVarshitha2025, TsaiLo2025, Duvuru2026}. Therefore, by providing a structured comparison of selected models and engineered features, this study aims to address the limitation of previous work that often focuses mainly on predictive accuracy without fully considering feature influence and performance on unseen weather sensor data \parencite{Eren2025, Ma2025}.

\section{2.1 Machine Learning in Solar Power Generation Forecasting}

Solar power forecasting is a critical component of energy management, especially as renewable energy becomes increasingly integrated into electrical grids. It involves predicting the amount of power generated by photovoltaic systems based on meteorological variables such as solar irradiance, temperature, wind speed, humidity, pressure, and cloud cover. These variables introduce variability into solar generation because solar output changes according to daily solar cycles, seasonal conditions, and short-term atmospheric changes. Solar radiation, sunshine duration, Julian day, cloud cover, and temporal meteorological features have been identified as important inputs for photovoltaic generation and solar irradiance prediction \parencite{CarvalhoJunior2025, Pongphaw2025, TsaiLo2025}. Machine learning has become a promising tool for this purpose because it can model complex and nonlinear relationships between weather sensor data and solar power output \parencite{Eren2025, JnanaVarshitha2025, Alhassan2025}.

Among the commonly used ML models for solar forecasting are Long Short-Term Memory (LSTM) networks, Random Forest (RF), Gradient Boosting Machines (GBM), XGBoost, LightGBM, CatBoost, and Support Vector Regression (SVR). LSTM, as a recurrent neural network architecture, is well suited for time-series forecasting because it can learn sequential dependencies from historical weather and solar generation data \parencite{Eren2025, TsaiLo2025, Duvuru2026}. In contrast, ensemble models such as Random Forest, XGBoost, LightGBM, and CatBoost are effective for structured tabular data because they can capture nonlinear feature interactions between multiple weather variables \parencite{JnanaVarshitha2025, Bakht2025, Alhassan2025}.

In addition to model selection, interpretability has become an increasingly important requirement in solar power forecasting. While complex models such as LSTM and gradient boosting models can produce accurate predictions, they are often considered black-box models because their decision-making process is not directly transparent. Explainable Artificial Intelligence techniques, such as SHAP and LIME, are used to improve transparency by showing how input features contribute to model predictions. SHAP- and LIME-based forecasting studies have used these methods to identify influential weather and temporal variables, explain individual predictions, and improve trust in model behaviour \parencite{JnanaVarshitha2025, Eren2025, TsaiLo2025, Duvuru2026}. In addition, meteorological feature contribution studies show that variables such as cloud-related features, humidity, and dew point can influence solar-related prediction tasks \parencite{Pongphaw2025}.

Furthermore, transfer learning and cross-station modelling have been explored to improve the generalizability of solar forecasting models across different geographical locations. This approach allows models trained on one dataset or meteorological station to be adapted to another location with limited additional data, which is useful when local weather sensor data are insufficient. For example, LSTM-based transfer learning has been used to improve short-term solar irradiance prediction across different meteorological stations, showing that model generalizability can be improved when spatial and temporal variations are considered \parencite{Eren2025}. Therefore, robust model comparison, interpretability analysis, and careful evaluation on unseen data are important components of a reliable solar forecasting framework.

In conclusion, machine learning techniques, especially LSTM and ensemble methods such as XGBoost, LightGBM, CatBoost, and Random Forest, provide a strong framework for solar power forecasting. However, the literature also shows that model performance depends on dataset characteristics, preprocessing procedures, engineered features, and evaluation design. Therefore, accuracy alone is not sufficient to justify model selection. A reliable forecasting framework should also consider whether the selected model can generalize to unseen weather sensor data and whether its predictions can be interpreted through feature influence analysis \parencite{ArellanosTafur2025, Eren2025, Ma2025}. The current study will evaluate selected machine learning models using the same weather sensor dataset, preprocessing procedure, feature engineering process, and evaluation metrics, while also analysing the influence of key features on forecasting performance.

\section{2.2 Overview of Machine Learning Algorithms for Solar Power Generation Forecasting}

This section reviews the main machine learning algorithms used in solar power generation forecasting, including Support Vector Regression, Random Forest, Gradient Boosting models such as XGBoost, LightGBM, CatBoost, and Long Short-Term Memory networks. These models are selected because they represent different modelling approaches: kernel-based regression, tree-based ensemble learning, boosting-based optimisation, and deep sequential learning. Ensemble and boosting models are commonly used for structured weather-based forecasting because they can capture nonlinear feature interactions, while LSTM models are relevant for time-series forecasting because they can learn sequential dependencies from historical weather and solar generation data \parencite{JnanaVarshitha2025, Bakht2025, Alhassan2025, Eren2025, TsaiLo2025, Duvuru2026}.

The main argument in this section is that algorithm selection should be based on the forecasting problem rather than model popularity. A model that performs well in one study may not perform equally well in another because solar forecasting datasets differ in weather variables, temporal resolution, geographical context, sensor quality, preprocessing method, feature construction, forecasting horizon, and evaluation metrics. For example, LightGBM may perform strongly on structured tabular forecasting data, while LSTM may be more suitable when the dataset contains strong sequential dependencies. Arellanos-Tafur et al. also showed that statistical approaches such as SARIMA can provide competitive forecasting performance and uncertainty quantification, which reinforces the point that deep learning should not be assumed to be superior in every forecasting context. Therefore, this study compares selected models under the same dataset and evaluation procedure to ensure that the performance comparison is fair and meaningful \parencite{ArellanosTafur2025, TsaiLo2025, CarvalhoJunior2025}.

Traditional machine learning models such as SVR and Random Forest are useful because they provide baseline performance with relatively lower complexity. SVR can model nonlinear regression relationships through kernel functions, while Random Forest can handle nonlinear interactions and provide feature importance. These characteristics are useful when working with weather sensor variables, especially when the dataset is moderate in size. However, these models do not naturally understand temporal sequence unless lag features, rolling averages, or time-based variables are explicitly included during feature engineering \parencite{JnanaVarshitha2025, Alhassan2025, CarvalhoJunior2025}.

Boosting models such as XGBoost, LightGBM, and CatBoost have become popular because they often achieve high accuracy on structured data. They are effective at learning nonlinear patterns and interactions among weather variables, making them suitable for weather-sensor-based solar forecasting tasks \parencite{JnanaVarshitha2025, Bakht2025, Alhassan2025}. However, like Random Forest, these models are not automatically time-series models. Their performance depends heavily on the quality of feature engineering, especially when the solar forecasting task involves daily cycles, seasonal patterns, or short-term weather persistence. Therefore, the value of boosting models in this study depends not only on their prediction accuracy but also on whether their predictions can be interpreted using XAI methods such as SHAP or LIME \parencite{TsaiLo2025, Pongphaw2025, JnanaVarshitha2025}.

Deep learning models such as LSTM are different because they are designed to learn sequential dependencies. This is highly relevant to solar forecasting because PV generation follows time-based patterns such as sunrise increase, midday peak, evening decline, seasonal variation, and weather transitions. However, LSTM usually requires more data, tuning, and computational resources than traditional models. This creates a trade-off between temporal modelling power and practical implementation complexity \parencite{Eren2025, TsaiLo2025, Duvuru2026}.

Overall, the literature suggests that the most defensible approach is not to claim one algorithm as universally superior, but to compare algorithms based on accuracy, robustness, interpretability, and suitability for weather sensor data. This directly aligns with the second objective of this study, which is to compare selected machine learning models for solar power generation forecasting using the same dataset, preprocessing procedure, feature engineering process, and evaluation metrics. A strong comparison will help determine whether simpler models such as SVR and Random Forest are sufficient, whether boosting models provide meaningful improvement, or whether sequential models such as LSTM are more suitable for capturing temporal patterns in the selected dataset \parencite{ArellanosTafur2025, Bakht2025, JnanaVarshitha2025, Alhassan2025, Eren2025}.

\subsection{2.2.1 Support Vector Regression (SVR)}

Support Vector Regression (SVR) is a supervised learning algorithm that can handle nonlinear regression problems, making it useful for predicting solar power from weather data. SVR works by finding a function that fits the data within a defined margin of tolerance, while kernel functions allow it to capture nonlinear relationships between input variables and output values. In solar forecasting, SVR can be useful when the dataset is moderate in size and when the relationship between weather variables and solar output can be represented through engineered features \parencite{JnanaVarshitha2025, Alhassan2025}.

While SVR can perform well in smaller datasets with relatively structured relationships between input variables and solar power, it may struggle when the data contain strong temporal dependencies. Solar power forecasting often requires the model to capture daily and seasonal changes in solar irradiance, as well as short-term weather transitions. Since SVR does not inherently model temporal order, features such as lagged values, rolling averages, hour of day, and day of year must be added to represent time-dependent behaviour. Previous studies on temporal encoding, lagged signals, and LSTM-based forecasting show that time-related information is important for representing solar forecasting patterns \parencite{TsaiLo2025, Eren2025, JnanaVarshitha2025}.

In terms of interpretability, SVR offers a simpler modelling structure compared with deep learning models, but it is less transparent than fully interpretable statistical models. In many comparative studies, SVR is used as a baseline model to evaluate whether more advanced models provide significant improvement. This is important because complex models such as LSTM or gradient boosting models may not always provide better performance if the dataset is small, noisy, or sufficiently represented by engineered features \parencite{ArellanosTafur2025, JnanaVarshitha2025, Alhassan2025}. Therefore, SVR remains useful as a benchmark for evaluating model complexity and forecasting performance.

For this study, SVR supports the objective of comparing model performance under the same weather sensor dataset. Its role is not necessarily to outperform all models, but to provide a controlled baseline for evaluating whether advanced models genuinely improve performance or only add complexity. If SVR performs competitively, it may indicate that the engineered weather-based and time-based features are already strong enough to represent the forecasting problem. If it performs poorly, it may suggest that more flexible nonlinear or sequential models are required \parencite{JnanaVarshitha2025, Alhassan2025, Bakht2025}.

\subsection{2.2.2 Random Forest and Decision Tree-Based Models}

Random Forest (RF) is an ensemble learning method that combines multiple decision trees to improve prediction accuracy and reduce overfitting. RF models are particularly strong in handling noisy datasets, nonlinear relationships, and feature interactions, making them a popular choice for solar power forecasting. RF can also provide feature importance scores, which are useful for identifying the relative influence of weather variables such as solar irradiance, temperature, humidity, wind speed, and cloud-related features \parencite{JnanaVarshitha2025, Bakht2025, Alhassan2025, Pongphaw2025}.

However, Random Forest models have limitations when capturing temporal dependencies in time-series data. While RF can be used for solar power forecasting by incorporating features such as lagged values, rolling statistics, and time-of-day indicators, it does not inherently learn sequential patterns in the same way as LSTM. For example, seasonal variation and daily solar cycles may not be effectively represented unless temporal features are explicitly engineered. This limitation is important because solar forecasting depends not only on current weather readings, but also on recent weather patterns and time-based behaviour \parencite{Eren2025, TsaiLo2025, CarvalhoJunior2025}.

Despite these limitations, Random Forest remains a strong candidate for solar forecasting because it offers a balance between predictive performance and interpretability. It can handle multiple weather sensor variables and identify influential predictors without requiring strict assumptions about the data distribution. For more complex forecasting tasks, however, models such as XGBoost, LightGBM, CatBoost, or LSTM may provide better performance if the dataset contains strong nonlinear or temporal patterns. The current study will compare Random Forest against other selected models to assess its feasibility for weather-sensor-based solar power forecasting \parencite{JnanaVarshitha2025, Bakht2025, Pongphaw2025, Alhassan2025}.

\subsection{2.2.3 Gradient Boosting Models: XGBoost, LightGBM, and CatBoost}

Gradient Boosting Models (GBMs), including XGBoost, LightGBM, and CatBoost, are powerful tools for solar power forecasting because they can model complex nonlinear relationships and handle structured weather sensor datasets effectively. These models use a sequential learning approach, where each new tree attempts to correct the errors of previous trees. This allows boosting models to capture interactions among weather variables and improve prediction accuracy in forecasting tasks \parencite{JnanaVarshitha2025, Bakht2025, Alhassan2025}.

XGBoost is widely used because of its strong predictive performance, regularisation capability, and robustness against overfitting. LightGBM is known for its efficient performance on structured forecasting data, while CatBoost can also model complex feature interactions in tabular datasets. In solar forecasting studies, these models are often compared with SVR, Random Forest, neural networks, and statistical models to determine which approach provides the best performance for a specific dataset \parencite{JnanaVarshitha2025, Bakht2025, Alhassan2025, CarvalhoJunior2025}.

Despite their strong predictive performance, GBMs such as XGBoost and LightGBM are often viewed as black-box models because their internal decision process is not directly obvious. To address this issue, SHAP and LIME can be used to improve model transparency and provide insights into how different weather features influence predictions. These techniques help identify whether variables such as solar irradiance, temperature, humidity, wind speed, or lagged features are important to the model. Therefore, gradient boosting models are relevant to this study not only because of their accuracy, but also because their predictions can be interpreted using XAI methods \parencite{JnanaVarshitha2025, TsaiLo2025, Pongphaw2025, Bakht2025}.

\subsection{2.2.4 Long Short-Term Memory (LSTM) Networks}

Long Short-Term Memory (LSTM) networks are specifically designed for sequential data, making them suitable for time-series forecasting tasks such as solar power generation prediction. LSTM models can learn temporal dependencies from historical weather sensor data, which is important because solar power output is affected by daily, seasonal, and weather-driven variations. This ability to retain information across time steps makes LSTM useful for forecasting problems where previous weather conditions influence future solar output \parencite{Eren2025, TsaiLo2025, Duvuru2026}.

Compared with traditional machine learning models, LSTM can learn temporal relationships more directly. Models such as SVR, Random Forest, and XGBoost require manually engineered lag features to represent previous weather conditions. LSTM, by contrast, can learn from ordered input sequences and retain useful information across time steps. This is important when solar forecasting depends on recent patterns, such as previous irradiance readings, previous cloud conditions, recent wind speed variation, or cyclic time-based patterns \parencite{Eren2025, TsaiLo2025, Duvuru2026}.

Recent studies support the relevance of LSTM in solar irradiance and PV forecasting. Eren et al. used LSTM-based models with transfer learning and SHAP to improve short-term solar irradiance forecasting across meteorological stations. Tsai and Lo developed an LSTM-SHAP framework using multivariate meteorological data, including temperature, humidity, pressure, wind conditions, global radiation, and temporal encodings. Duvuru et al. applied a multi-layer LSTM model with SHAP-based explainability for PV cell power forecasting using Global Horizontal Irradiance data. These studies show that LSTM is useful when temporal dynamics are central to the forecasting problem \parencite{Eren2025, TsaiLo2025, Duvuru2026}.

However, LSTM also has weaknesses. It usually requires more data, more training time, more tuning, and higher computational resources than traditional models. If the dataset is small or noisy, LSTM may overfit or fail to outperform simpler methods. Therefore, it is not appropriate to assume that LSTM is automatically the best model for all solar forecasting tasks. Arellanos-Tafur et al. showed that SARIMA could provide strong uncertainty quantification and competitive accuracy compared with some deep learning alternatives. This reinforces the need for careful model comparison rather than model bias \parencite{ArellanosTafur2025, CarvalhoJunior2025, Eren2025}.

Interpretability is another challenge for LSTM. Because LSTM learns internal sequential representations, it can be difficult to explain why it produces a specific prediction. This is problematic in solar forecasting because users may need to understand whether the forecast is mainly influenced by irradiance, humidity, temperature, wind speed, or time-based factors. SHAP has been used with LSTM-based forecasting to identify influential variables, such as lagged wind speed, cyclic time-based features, temperature, humidity, and other meteorological inputs. Therefore, explainability is necessary if LSTM is used in an operationally meaningful forecasting framework \parencite{Eren2025, TsaiLo2025, Duvuru2026}.

\section{2.3 Feature Engineering in Solar Power Forecasting}

Feature engineering plays a critical role in solar power forecasting because raw sensor data, such as solar irradiance, temperature, humidity, wind speed, and pressure, may not fully capture the complex relationships that affect solar power generation. Before feature engineering can be performed, weather sensor data must be preprocessed to address missing values, noisy readings, outliers, duplicated records, and inconsistent timestamps, because poor data quality can lead to unreliable forecasting performance. After preprocessing, weather sensor variables can be transformed into meaningful weather-based and time-based features that allow machine learning models to recognise daily cycles, seasonal variation, short-term persistence, and nonlinear interactions. Previous solar forecasting studies support the importance of structured meteorological inputs, temporal features, lagged signals, and explainable feature analysis in improving forecasting performance and interpretation \parencite{TsaiLo2025, JnanaVarshitha2025, Duvuru2026}.

Previous studies emphasise the importance of various weather features in predicting solar power. Carvalho Junior et al. demonstrated that solar radiation, sunshine duration, and Julian day are important variables for PV generation prediction. Tsai and Lo showed that temporal features, such as hour of day and day of year, help LSTM-based models capture diurnal and seasonal patterns in solar irradiance. Similarly, Jnana Varshitha et al. highlighted that lagged signals can contribute to forecasting performance and model interpretation. These studies support the idea that feature engineering transforms raw weather sensor readings into usable inputs that improve predictive accuracy and feature influence analysis \parencite{CarvalhoJunior2025, TsaiLo2025, JnanaVarshitha2025}.

Feature engineering in solar forecasting also requires the creation of temporal features that capture predictable cycles in solar energy production. Solar power generation exhibits daily and seasonal patterns, with output usually increasing after sunrise, peaking around midday, and declining toward evening. These cyclical patterns can be represented using features such as hour of day, day of year, month, seasonal indicators, and cyclic encodings. The inclusion of lagged features, such as previous irradiance, previous temperature, or previous wind speed values, helps models capture short-term persistence effects and anticipate future changes based on recent weather conditions. By incorporating these features, models can adapt better to daily and seasonal variation \parencite{TsaiLo2025, JnanaVarshitha2025, Eren2025}.

The role of additional weather variables, such as cloud cover, temperature, humidity, dew point, and wind speed, is also important in improving forecasting accuracy. Pongphaw et al. demonstrated that cloud-related and humidity-related meteorological variables can provide useful feature contribution in cloud cover prediction, while Carvalho Junior et al. showed that sunshine duration and solar radiation are strongly associated with photovoltaic energy generation. Similarly, SHAP-based solar forecasting studies have shown that meteorological variables such as humidity, wind speed, lagged signals, and cyclic time-based features can influence forecasting outputs. Therefore, using weather variables alongside solar irradiance improves model robustness by accounting for both direct and indirect factors affecting solar generation \parencite{Pongphaw2025, CarvalhoJunior2025, Eren2025, TsaiLo2025, Duvuru2026}.

For this study, feature engineering will be a core focus in improving forecasting accuracy and interpretability. This study will explore the impact of different combinations of weather sensor data and engineered features, such as lagged values, rolling averages, and time-based encodings, on model performance. The study will also investigate the importance of key weather variables such as solar irradiance, temperature, humidity, wind speed, pressure, and cloud-related features in forecasting solar power generation. By testing how these features influence model predictions, the research will address the gap in the literature regarding systematic feature engineering and feature influence analysis. This directly supports the first research objective of preprocessing and engineering suitable weather-based and time-based features for solar power generation forecasting \parencite{CarvalhoJunior2025, TsaiLo2025, JnanaVarshitha2025, Pongphaw2025}.



\section{2.4 Explainable Artificial Intelligence (XAI) in Solar Power Forecasting}

Explainable Artificial Intelligence (XAI) has become a crucial area in machine learning, especially for complex forecasting tasks such as solar power prediction. While models such as LSTM, XGBoost, LightGBM, and CatBoost can provide high prediction accuracy, their black-box nature can make it difficult for users to understand how forecasts are produced. This lack of interpretability presents a challenge in solar power forecasting because grid operators, energy managers, and system planners may need to understand which weather variables are influencing the model's prediction. XAI methods such as SHAP and LIME are increasingly used to improve transparency, trust, and decision-making in ML-based energy forecasting \parencite{JnanaVarshitha2025, Eren2025, TsaiLo2025, Duvuru2026}.

SHAP is one of the most widely used techniques in XAI because it assigns a contribution value to each feature, indicating how much that feature contributes to a model prediction. In solar power forecasting, SHAP can help explain the influence of individual weather variables such as solar irradiance, temperature, humidity, cloud cover, wind speed, and time-based features on forecasted output. For example, Eren et al. used SHAP to identify influential lagged wind speed and cyclic time-based features in LSTM-based solar irradiance forecasting. Tsai and Lo also used SHAP to interpret an LSTM-based solar irradiance prediction framework using multivariate meteorological inputs, while Duvuru et al. applied SHAP-based explainability in photovoltaic power forecasting using Global Horizontal Irradiance data \parencite{Eren2025, TsaiLo2025, Duvuru2026}.

LIME is another XAI method that approximates complex machine learning models with simpler local models to explain individual predictions. This local approximation helps users understand why a model produces a specific prediction for a specific data instance. For example, LIME can explain why a model predicts lower solar power generation under conditions such as low irradiance, high cloud cover, or unfavourable weather sensor readings. Jnana Varshitha et al. integrated both SHAP and LIME in an energy forecasting framework, where SHAP was used to quantify global and local feature contributions while LIME provided case-specific explanations for individual forecasts. This shows that LIME can complement SHAP when instance-level explanation is required \parencite{JnanaVarshitha2025}.

The integration of XAI techniques such as SHAP and LIME is not only about improving model interpretability; it can also support feature selection and model refinement. By identifying which features have the strongest influence on model predictions, researchers can focus on the most important weather variables and reduce unnecessary model complexity. For instance, Pongphaw et al. used SHAP to identify stable meteorological feature contributions in cloud cover prediction, particularly for dew point and relative humidity. Jnana Varshitha et al. also showed that lagged signals were important in energy forecasting and used explainability methods to support feature interpretation. These findings suggest that XAI can be used to understand model behaviour and guide feature selection, although the results should be interpreted as model-based feature influence rather than direct causal proof \parencite{Pongphaw2025, JnanaVarshitha2025, Eren2025}.

Moreover, XAI techniques are particularly useful for addressing interpretability challenges associated with deep learning models such as LSTM. Because LSTM models learn internal sequential representations, it is not always clear which input variables or time steps are most influential. By using SHAP, researchers can analyse whether predictions are driven by lagged irradiance, wind speed, temperature, humidity, pressure, or time-based encodings. This level of transparency is important for improving trust in forecasting models and for supporting practical solar energy applications. However, XAI results should be treated carefully because they explain how the trained model uses features within a dataset, not necessarily the physical cause of changes in solar output \parencite{Eren2025, TsaiLo2025, Duvuru2026}.

In conclusion, XAI techniques such as SHAP and LIME are essential for improving the interpretability and trustworthiness of solar power forecasting models. They allow researchers and users to understand how weather variables, such as solar irradiance, temperature, humidity, wind speed, and cloud-related features, contribute to model predictions. Furthermore, XAI supports feature selection and model refinement by identifying important features and reducing unnecessary complexity. For this study, the integration of XAI directly supports the third objective, which is to identify the best-performing forecasting model and analyse the influence of key weather sensor features on solar power generation prediction \parencite{JnanaVarshitha2025, Eren2025, TsaiLo2025, Duvuru2026}.




\subsection{2.5 Dataset Diversity and Model Generalizability}

Dataset diversity and model generalizability are important aspects of machine learning model evaluation in solar power forecasting. Solar power generation is affected by weather conditions, geographical location, sensor quality, temporal resolution, photovoltaic system characteristics, and forecasting horizon. As a result, a model that performs well on one dataset may not produce the same level of accuracy when applied to another dataset or under different environmental conditions. This shows that solar forecasting models should not be judged only based on high accuracy in one experimental setting, but also on whether they can maintain reliable performance when tested on unseen data. Previous studies on cross-station learning, transfer learning, and robust photovoltaic forecasting further support the importance of evaluating generalization rather than relying only on training or validation accuracy \parencite{Eren2025, Ma2025}.

Previous studies show that forecasting performance varies across datasets and operating conditions. For example, Jnana Varshitha et al. showed that different machine learning models may perform differently depending on the forecasting dataset, which supports the need for customised model selection. Eren et al. demonstrated that cross-station learning and transfer learning can improve the generalization performance of LSTM-based solar irradiance forecasting. Similarly, Ma et al. showed that a gray-box photovoltaic model improved prediction accuracy, robustness, interpretability, and generalization compared with several deep learning models. These findings suggest that model generalizability is a key issue in solar forecasting and should be considered together with prediction accuracy \parencite{JnanaVarshitha2025, Eren2025, Ma2025}.

Dataset diversity also affects the interpretation of model performance. Some studies use solar power generation data, while others use solar irradiance, photovoltaic cell power, or climatic station data. The input variables may also differ across studies, including solar irradiance, temperature, humidity, wind speed, pressure, cloud cover, sunshine duration, Julian day, or lagged signals. Because of these differences, it is not always fair to directly compare reported accuracy values from different papers. A model may appear better because it was tested on a cleaner dataset, a shorter forecasting horizon, or a simpler prediction task. Therefore, this study compares selected machine learning models using the same dataset, preprocessing procedure, feature engineering process, and evaluation metrics to produce a fairer model comparison \parencite{ArellanosTafur2025, CarvalhoJunior2025, Bakht2025}.

Generalizability is also closely related to data splitting and leakage prevention. In time-series forecasting, random data splitting may create overly optimistic results because future patterns can accidentally influence the training process. To reduce this risk, chronological splitting is more suitable because it preserves the time order of solar power generation data. Training data should be used to develop the model, while unseen testing data should be used to evaluate whether the trained model can maintain reliable performance beyond the training period. In addition, scaling, preprocessing, feature selection, and hyperparameter tuning should be fitted or decided using only the training and validation data before final testing. This helps prevent data leakage and supports a more realistic evaluation of forecasting performance \parencite{Eren2025, TsaiLo2025, Duvuru2026}.

Although external generalizability across different locations and photovoltaic systems is important, it is beyond the scope of this study. This research does not claim that the final model will generalize to all geographical regions, solar farms, weather sensor configurations, or photovoltaic technologies. Instead, generalizability in this study refers to the ability of the best-performing model to maintain reliable forecasting performance on unseen weather sensor data from the same dataset. This scope is appropriate for evaluating whether the selected model is overfitting to the training data or whether it can produce stable predictions on new data within the same experimental context.

In conclusion, dataset diversity and model generalizability are important because solar forecasting performance depends on dataset characteristics, environmental conditions, preprocessing procedures, feature engineering, model selection, and evaluation design. The literature shows that no single model can be assumed to be universally superior across all solar forecasting tasks. Therefore, this study addresses this issue by evaluating selected machine learning models under a consistent experimental framework and by analysing the best-performing model using unseen weather sensor data. This directly supports the third research objective, which is to identify the best-performing forecasting model and analyse its generalizability on unseen weather sensor data and the influence of key weather sensor features on solar power generation prediction \parencite{ArellanosTafur2025, Eren2025, Ma2025}.